% 05_particles.tex - save MODFLOW's flows and track particles backwards from the well
\section{Track particles from the pumping well}
MODFLOW~6's particle tracking (PRT) runs after the coupled run, on the flows that MODFLOW saved.

\step{Save MODFLOW's flows every day and run}
Still in \cmd{\$WS/work} (with the ensemble lines removed), run one year and ask Raven to save MODFLOW's flows every day.
Change \cmd{Creek.rvi} and \cmd{Creek.rvg}:
\diff{t5_prt_rvi.diff}
\diff{t5_prt.diff}
\term{t5_run_prt.txt}
\see \cmd{Exiting Gracefully: Successful Simulation}.
\why \cmd{out\_prt/mf6/gwf.cbc} (the \emph{budget file}) now holds the flows of every package, the specific discharge and
the saturation of every day (\V{t5:cbcmb}~MB for this year). \cmd{gwf.porosity} holds the porosity of every cell
(\cmd{POROSITY} of the aquifer classes).

\step{Track particles backwards}
\cmd{make\_prt.py} writes a PRT simulation for these flows and runs it with the library. By default it releases
\V{t5:npart} particles around each well of \cmd{:Wells}. \cmd{--backward} tracks them backwards in time, to where the
pumped water came from:
\term{t5_make_prt.txt}
\term[basicstyle=\ttfamily\tiny]{t5_endpoints.txt}
\why each line is a particle: when it was released (in days of the tracking), its release point and its end point, the
time it was tracked (here the whole year: no particle left the model), the PRT status code of its last record and, in
words, why its track ended. The points are in metres, in the frame of the grid file \cmd{gwf.dis.grb} (PRT's own
coordinates shifted by its \cmd{XORIGIN} and \cmd{YORIGIN}).
\begin{trouble}
\begin{itemize}
\item \emph{``No module named flopy''}: type \cmd{pip install flopy} with the virtual environment active.
\item \cmd{make\_prt.py} stops with a message when the flows were not saved every day: add \cmd{:HeadSaveFrequency 1} and
  \cmd{:WriteGWBudgetFile} to the \cmd{.rvg}.
\item PRT needs a regular or a quadtree grid (MODFLOW's DIS or DISV). HRU meshes and imported grids use DISU, which PRT
  cannot track.
\end{itemize}
\end{trouble}

\step{Plot the tracks}
The script \cmd{plot\_tracks.py} draws the tracks, in the same frame, with the travel times of \cmd{endpoints.csv}:
\term{t5_plot_tracks.txt}
\fig[0.62]{t5_tracks.png}{\cmd{out\_prt/prt/tracks.png}: backward tracks from the \V{t5:npart} release points around PW\_1
over one year. Going back over the year, each particle moves \V{t5:distmin} to \V{t5:distmax}~m from its release point.}
\begin{note}[Tracks need small cells]
All the tracks begin and end inside the well's own 250-m cell. PRT interpolates the flow within a cell, and this cell is
itself the sink of the well. The tracks show the method; they are not a capture zone. A study of capture zones or travel
times needs cells much smaller than the distances of interest around the well (\cmd{:GridRefinement WELLS}, or a smaller
\cmd{:GridCellSize}).
\end{note}
